\documentclass[11pt]{article}

\usepackage[margin=1in]{geometry}
\usepackage{amsmath,amssymb}
\usepackage{graphicx}
\usepackage{booktabs}
\usepackage{siunitx}
\usepackage{xcolor}
\usepackage{tcolorbox}
\usepackage[hidelinks]{hyperref}
\usepackage{enumitem}

\newcommand{\Fk}{F^{k}}
\newcommand{\Gk}{G^{k}}
\newcommand{\Rk}{R^{k}}
\newcommand{\znl}{\zeta_{n\ell}}
\newcommand{\Eav}{E_{\mathrm{av}}}
\newcommand{\prog}[1]{\textsc{#1}}
\newcommand{\file}[1]{\texttt{#1}}
% spectroscopic ion notation, e.g. \ion{Fe}{2} -> "Fe II"
\newcommand{\ion}[2]{\mbox{#1\,\textsc{\romannumeral #2}}}

\newtcolorbox{keybox}[1]{colback=blue!4,colframe=blue!40!black,title=#1,
  fonttitle=\bfseries,boxrule=0.4pt}

\title{\textbf{The Cowan Atomic-Structure Codes}\\[2pt]
\large What \prog{rcn}, \prog{rcn2}, \prog{rcg}, \prog{rce} Actually Compute,
from First Principles\\[6pt]
\normalsize A working explainer for the Modern Atomic Linelists (\texttt{mal}) project}
\author{C.\ Conroy \&\ Claude}
\date{\today}

\begin{document}
\maketitle

\begin{abstract}
\noindent
This document explains, from first principles, what the four Cowan programs do
and how they chain together to produce atomic energy levels and oscillator
strengths ($gf$-values) --- the data that go into stellar line lists. It maps the
physics onto the \emph{actual} programs, input decks, and output files in this
repository (\file{cowan\_nist/extracted/}). The goal is a shared, durable
understanding of the engine we are modernizing, with particular attention to the
least-squares fitting stage (\prog{rce}), which is the heart of the project. The
canonical reference throughout is R.~D.\ Cowan, \emph{The Theory of Atomic
Structure and Spectra} (Univ.\ California Press, 1981), cited as \textbf{TASS}.
\end{abstract}

\tableofcontents

\section{The big picture: why these four programs exist}

A stellar line list needs, for each transition, the lower and upper energy
levels, the wavelength, and the oscillator strength $gf$. The wavelength is just
the level difference; the levels and the $gf$ are the hard part. Computing them
\emph{ab initio} for complex atoms (open $d$- and $f$-shells, the iron group) is
not reliably accurate. Cowan's insight, embodied in these codes, is a
\emph{semi-empirical} scheme: compute the atomic structure approximately from
theory, then \emph{adjust a small number of physically meaningful radial
parameters until the predicted levels match the experimentally measured ones}.
The wavefunctions (eigenvectors) that result are far better than the purely
theoretical ones, and the $gf$-values computed from them are correspondingly
better.

\begin{keybox}{The pipeline in one breath}
\prog{rcn} solves for the radial wavefunctions and the \emph{radial integrals}
($\Fk,\Gk,\znl,\dots$) of each electron configuration $\;\to\;$ \prog{rcn2}
forms the configuration-interaction and electric-multipole (transition) radial
integrals between configurations $\;\to\;$ \prog{rcg} builds the angular part,
assembles and \emph{diagonalizes} the energy (Hamiltonian) matrix to get energy
levels and eigenvectors, and from the eigenvectors computes $gf$-values
$\;\to\;$ \prog{rce} re-runs the diagonalization inside a \emph{least-squares
loop}, varying the radial parameters until the computed levels best fit the
observed levels.
\end{keybox}

The division of labor reflects a deep structural fact about atomic Hamiltonians:
the energy of a level separates into \textbf{radial} factors (integrals over the
one-electron radial wavefunctions --- these are the $\Fk$, $\Gk$, $\znl$,
$\Rk$) multiplied by \textbf{angular} factors (pure angular-momentum algebra,
the same for any atom with a given configuration). \prog{rcn}/\prog{rcn2}
compute the radial side; \prog{rcg} computes the angular side and combines them;
\prog{rce} treats the radial numbers as adjustable.

\section{The physics each stage implements}

\subsection{The central-field starting point and \prog{rcn}}

We seek the eigenstates of the non-relativistic many-electron Hamiltonian
\begin{equation}
H = \sum_{i}\left(-\tfrac{1}{2}\nabla_i^2 - \frac{Z}{r_i}\right)
  + \sum_{i<j}\frac{1}{r_{ij}}
  + \sum_i \xi(r_i)\,(\mathbf{l}_i\!\cdot\!\mathbf{s}_i),
\label{eq:H}
\end{equation}
in atomic units, where the last term is the spin--orbit interaction. The
\emph{central-field approximation} writes each electron's wavefunction as a
product of a radial part $P_{n\ell}(r)/r$ and angular/spin parts. \prog{rcn}
solves the resulting one-electron radial equations self-consistently for a
\emph{single configuration} (e.g.\ $3d^7$, or $3d^6 4p$), using one of several
approximations to Hartree--Fock; in practice the Hartree-plus-statistical-%
exchange (HX) method or center-of-gravity HF is used, optionally with
relativistic corrections (the ``HFR'' / ``HXR'' options --- important for the
iron group and heavier).

From the converged radial wavefunctions $P_{n\ell}(r)$, \prog{rcn} computes the
\textbf{radial integrals} that parametrize the energy:
\begin{align}
\Fk(n\ell,n'\ell') &= \int\!\!\int P_{n\ell}^2(r_1)\,
  \frac{r_<^{k}}{r_>^{k+1}}\,P_{n'\ell'}^2(r_2)\;dr_1\,dr_2
  &&\text{(direct Slater integral)},\\
\Gk(n\ell,n'\ell') &= \int\!\!\int P_{n\ell}(r_1)P_{n'\ell'}(r_1)\,
  \frac{r_<^{k}}{r_>^{k+1}}\,P_{n\ell}(r_2)P_{n'\ell'}(r_2)\;dr_1\,dr_2
  &&\text{(exchange Slater integral)},\\
\znl &= \tfrac{\alpha^2}{2}\!\int P_{n\ell}^2(r)\,
  \frac{1}{r}\frac{dV}{dr}\;dr
  &&\text{(spin--orbit parameter)},
\end{align}
plus $\Eav$, the configuration's center-of-gravity energy. These few numbers
per configuration are the \emph{radial parameters}. Everything semi-empirical
happens by treating them as adjustable later (\S\ref{sec:rce}).

\subsection{Configuration interaction and transition integrals: \prog{rcn2}}

Real atomic states are not pure single configurations; nearby configurations of
the same parity and $J$ \emph{mix} (configuration interaction, CI). \prog{rcn2}
takes the radial wavefunctions from \prog{rcn} for \emph{several} configurations
and computes:
\begin{itemize}[nosep]
\item the \textbf{CI Slater integrals} $\Rk$ (the $r_{ij}^{-1}$ matrix elements
\emph{between} different configurations), which set the strength of the mixing;
\item the \textbf{electric-multipole radial integrals}
$\langle n\ell \,\|\, r \,\|\, n'\ell'\rangle$ (and higher multipoles), the
radial part of the transition dipole moment --- the seed of every $gf$-value.
\end{itemize}

\subsection{Angular algebra, diagonalization, and where $gf$ comes from: \prog{rcg}}
\label{sec:rcg}

\prog{rcg} supplies the \emph{angular} half. Using angular-momentum
(Racah/Wigner) algebra and tabulated coefficients of fractional parentage
(the binary ``CFP decks'' \file{FOR072/073/074} that a freshly built \prog{rcg}
generates once), it builds the Hamiltonian matrix in a chosen coupling basis.
Schematically, every matrix element of \eqref{eq:H} is
\begin{equation}
\langle \gamma' \,|\, H \,|\, \gamma\rangle
= \sum_{\text{integrals}}
\underbrace{c_{\gamma'\gamma}^{(\cdot)}}_{\text{angular coeff.\ (\prog{rcg})}}
\;\times\;
\underbrace{\{\Eav,\Fk,\Gk,\znl,\Rk\}}_{\text{radial integrals (\prog{rcn}/\prog{rcn2})}} .
\label{eq:matrixelement}
\end{equation}
\prog{rcg} then \textbf{diagonalizes} this matrix for each $J$. The eigenvalues
are the predicted \textbf{energy levels}; the eigenvectors give the
\textbf{composition} of each level (its percentage breakdown over basis states
--- e.g.\ ``$72\%\;{}^5\!D + 18\%\;{}^3\!F + \dots$'').

The oscillator strength follows from the eigenvectors. The line strength between
levels $\beta J$ and $\beta' J'$ is
\begin{equation}
S = \left| \langle \beta J \,\|\, \mathbf{P}^{(1)} \,\|\, \beta' J'\rangle \right|^2 ,
\qquad
gf = \frac{8\pi^2 m_e c}{3 h}\,\frac{\sigma}{g}\,S \;\propto\; \sigma\, S,
\end{equation}
where $\sigma$ is the transition wavenumber and the reduced matrix element of
the dipole operator $\mathbf{P}^{(1)}$ is a sum over the
\prog{rcn2} transition radial integrals weighted by \emph{both} levels'
eigenvector components. This is the crucial coupling that drives the whole
project: \textbf{many lines' $gf$-values depend on the same few radial
parameters through the shared eigenvectors.} A small change in the fit
redistributes oscillator strength among coupled lines.

\subsection{The semi-empirical least-squares fit: \prog{rce}}
\label{sec:rce}

This is the stage we are modernizing, so we treat it carefully. \prog{rcg}'s
predicted levels are only as good as the theoretical radial integrals --- and
those are systematically off (Hartree--Fock overestimates the Slater integrals
by $\sim$10--20\%). \prog{rce} fixes this empirically. Its job (TASS \S16-3 to
16-5):

\begin{enumerate}[nosep]
\item Treat the radial parameters $\mathbf{x}=\{\Eav,\Fk,\Gk,\znl,\Rk,\dots\}$
as \emph{free}. The angular coefficients in \eqref{eq:matrixelement} are fixed
(precomputed by \prog{rcg} and passed in the binary file \file{TAPE2E}).
\item For trial $\mathbf{x}$, build and diagonalize the energy matrix for each
$J$ to get computed levels $E_i^{\mathrm{calc}}(\mathbf{x})$.
\item Minimize the weighted least-squares residual against the
\emph{experimentally observed} levels $E_i^{\mathrm{obs}}$ (from NIST/lab),
\begin{equation}
\chi^2(\mathbf{x}) = \sum_i w_i \left[
E_i^{\mathrm{calc}}(\mathbf{x}) - E_i^{\mathrm{obs}} \right]^2 ,
\label{eq:chi2}
\end{equation}
iterating (Newton/Levenberg-style) until the parameters stop changing
(by $<0.03$ between cycles) or a cycle cap is hit. Parameters can be held fixed,
or tied in fixed ratios (a group forced to scale together).
\item Feed the \emph{fitted} $\mathbf{x}$ back through the diagonalization to get
the \emph{improved} eigenvectors --- and hence the improved $gf$-values
(via \prog{rcg}).
\end{enumerate}

\begin{keybox}{Why this is ``black magic'' --- and our opportunity}
The fit in \eqref{eq:chi2} is ill-conditioned: parameters are correlated, the
problem is non-linear, and which configurations to include and which parameters
to free/fix/tie is a matter of hand-tuned judgment with ``no strict
methodology.'' This is the only non-automated stage, takes hours to weeks per
ion, and has scarcely changed in 50 years. It is exactly where modern
techniques --- regularization, Bayesian inference with priors, MCMC for
\emph{per-line $gf$ uncertainties}, global optimization, and (our novel idea)
folding the observed \emph{stellar spectrum} into the objective alongside the
energy levels --- can make a genuine advance. See the project plan.
\end{keybox}

\section{Mapping to the actual code and files}

\begin{table}[h]
\centering
\small
\begin{tabular}{@{}llp{7.4cm}@{}}
\toprule
Program & Source (\file{for/}) & Computes \\
\midrule
\prog{rcn}  & \file{RCN36K.F} & Single-config.\ radial wavefns $P_{n\ell}(r)$;
  radial integrals $\Eav,\Fk,\Gk,\znl$ \\
\prog{rcn2} & \file{RCN2K.F}  & CI integrals $\Rk$; multipole transition
  integrals $\langle r\rangle$ \\
\prog{rcg}  & \file{rcg11k.f} & Angular coeffs (CFP); builds \& diagonalizes
  $H$ $\to$ levels, eigenvectors, $gf$ \\
\prog{rce}  & \file{RCE20K.F} & Least-squares fit of radial parameters to
  observed levels; improved eigenvectors $\to$ $gf$ \\
\bottomrule
\end{tabular}
\caption{The four programs and what they compute. Built by
\file{build/build.sh}.}
\end{table}

\paragraph{Data flow between programs.} Each program writes the input for the
next, mostly through fixed-name files (a legacy of the punched-card/tape era):
\begin{itemize}[nosep]
\item \prog{rcg} consumes the binary CFP decks \file{FOR072/073/074} (generated
once by \file{build/make\_cfp.sh}) and writes, among others, \file{TAPE2E}
(binary coefficient matrices for \prog{rce}) and \file{OUTGINE} (a starter for
\prog{rce}'s text input).
\item \prog{rce} reads \file{TAPE2E} plus a text deck \file{INE20} (control
info + the \emph{observed} levels --- the only thing not produced upstream) and
writes \file{LEVELS1/2/3} (fitted levels and eigenvectors) and \file{PARVALS}
(fitted parameter values, for restarting).
\end{itemize}

\paragraph{An actual input deck.} The example \file{work/IN36} is an \prog{rcn}
input for \ion{Sn}{7+} (\file{IN36.ai} is \ion{B}{3+} with a long Rydberg
series). The first line is a fixed-column ``control card'' (mesh, method flags,
convergence tolerances, relativistic/correlation switches); each subsequent line
names one configuration. For \ion{Sn}{7+}:
\begin{verbatim}
200-90 0 2  01.  4.0    5.E-08    1.E-11-2 00090 0 1.0  0.65  0.0 1.00   -6
   50    8Sn7+  4d7              4p6  4D7
   50    8Sn7+  4d65p            4p6  4D6  5P
   -1
\end{verbatim}
Here the two configurations $4d^7$ and $4d^6 5p$ (opposite parity --- the source
and upper levels of E1 transitions) are listed, terminated by \texttt{-1}. The
columns of this card are documented in \file{cowan\_nist/extracted/RCN\_DOC.txt};
we will annotate them in detail when we drive a real run.

\section{What ``running it'' looks like end to end}

For a single ion and one parity pair, the canonical chain is
\[
\prog{rcn} \to \prog{rcn2} \to \prog{rcg} \to \prog{rce} \to \prog{rcg}
\]
(the final \prog{rcg} re-run applies the fitted parameters to emit final levels
and $gf$). The seven coupling schemes \prog{rcg}/\prog{rce} support
(LS, JJ, intermediate, \dots) only change the basis, not the physics; final
eigenvectors are reported in LS or JJ for interpretation.

\begin{keybox}{Next steps in the project}
\begin{enumerate}[nosep]
\item Drive the shipped example (\file{IN36}) through the full chain and read
each output file together, annotating the real column formats.
\item Reproduce a known Kurucz \ion{Fe}{2} $gf$ via the standard chain (validate
we replicate the legacy result).
\item Replace the \prog{rce} fit \eqref{eq:chi2} with a modern regularized /
Bayesian fit, attach per-line $gf$ uncertainties, and (Tier~2) fold in the
solar spectrum.
\end{enumerate}
\end{keybox}

\section*{References}
\begin{itemize}[nosep]
\item R.~D.\ Cowan, \emph{The Theory of Atomic Structure and Spectra}
(Univ.\ California Press, 1981). [TASS]
\item Program documentation shipped with the code:
\file{RCN\_DOC.txt}, \file{RCG\_DOC.txt}, \file{RCE\_DOC.txt}
(in \file{cowan\_nist/extracted/}).
\item A.\ Kramida, ``Cowan Code: 50 Years of Growing Impact on Atomic Physics,''
\emph{Atoms} \textbf{7}, 64 (2019).
\end{itemize}

\end{document}
